% Generated from main.typ. Typst is authoritative; regeneration overwrites this file.
\documentclass{qlnotes-export}

\title{MATH 597: Measure Theory}
\subtitle{Typst-first course notes and worked homeworks}
\author{Qiulin Fan}
\date{Winter 2025}

\begin{document}
\mainmatter
\chapter{\texorpdfstring{the dual of \(L^{p}\) spaces}{the dual of L\^{}\{p\} spaces}}\label{the-dual-of-lp-spaces}

\section{\texorpdfstring{the dual of \(L^{p}\)-I {[}Fol 6.2{]}}{the dual of L\^{}\{p\}-I {[}Fol 6.2{]}}}

对应: Folland 5.1, 6.2.\\
(原本这是在 lec 25 的位置讲的, 但是当时由于没有 Radon-Nikodym Thm, 没有足够的工具去完成

\[(L^{p})^{\ast} = L^{q}\]

的证明 (差了一个 proof surjectivity of the isometry \(g\rightarrow\ell_{g}\)). 因而我把它放在这里, 衔接下面几个 lectures, 完成 6.2 这一节.\\
我们首先练习一个 example of Hölder's ineq 来回忆一下:\\
recall Hölder's ineq: for \(1 \leq p,q \leq \infty,\frac{1}{p} + \frac{1}{q} = 1\Longrightarrow\)

\[\left. | \middle| fg \middle| \middle| {}_{1} \leq \middle| \middle| f \middle| \middle| {}_{p} \middle| \middle| g \middle| |_{q} \right.\]

% qlnotes: kind=example; id=ex-13-the-dual-of-l-p-spaces-example-001; concepts=example-001
\begin{example}\label{ex-13-the-dual-of-l-p-spaces-example-001}

Prove:

\[f \in L^{3}(\lbrack - 1,1\rbrack,m)\Longrightarrow\int_{- 1}^{1}\frac{|f(x)|}{\sqrt{|x|}}\, dx < \infty\]

\end{example}

\begin{proof}

Apply Hölder's: 既然 \(f \in L^{3}\), 那么我们就拉满, take \(p = 3\), correspondingly \(q = 3/2\):

\[\left. \int_{- 1}^{1}\frac{|f(x)|}{\sqrt{|x|}}\, dx \leq (\int_{- 1}^{1} \middle| f(x) \middle| {}_{3}\, dx)^{\frac{1}{3}}(\int_{- 1}^{1}\frac{1}{|x|^{\frac{3}{4}}}\, dx)^{\frac{2}{3}} \right.\]

both integrals evaluate \(< \infty\)

\end{proof}

\subsection{intro to dual space}\label{intro-to-dual-space}

这里只讨论 \({\mathbb{K}} := {\mathbb{R}}\) or \(\mathbb{C}\).\\
recall, 对于一个 \(\mathbb{K}\)-vector space \(V\), 一个 linear functional of \(V\) 就是一个 linear function

\[f:V\rightarrow{\mathbb{K}}\]

对于作为 NVS 的 \(V\), 我们还可以定义一个 linear functional 的 boundedness.

% qlnotes: kind=definition; id=def-13-the-dual-of-l-p-spaces-bounded-linear-functional; concepts=bounded-linear-functional; aliases=bounded linear functional
\begin{definition}{bounded linear functional}\label{def-13-the-dual-of-l-p-spaces-bounded-linear-functional}

Let \(V\) be a \(\mathbb{K}\)-NVS, \(f:V\rightarrow{\mathbb{K}}\) be a linear functional.\\
我们称 \(f\) bounded, if exist \(C > 0\) s.t.

\[\left. |f(v) \middle| \leq C \middle| \middle| v \middle| \middle| ,\quad\forall\, v \in V \right.\]

\end{definition}

\begin{qlremark}{Remark}

注意, \textbf{linear functional 的 boundedness 和它作为函数的 boundedness 是不一样的概念.}\\
作为函数的 boundedness 表示函数值的有界性, 而\textbf{作为 linear map 的 boundedness (此处) 表示它的作用效果的 boundedness, 不会把一个 vector 放大太多倍.}

\end{qlremark}

% qlnotes: kind=proposition; id=prop-13-the-dual-of-l-p-spaces-linear-functional-bounded-iff-ctn-at-0; concepts=linear-functional-bounded-iff-ctn-at-0; aliases=linear functional bounded \iff ctn at 0
\begin{proposition}{linear functional bounded \(\Leftrightarrow\) ctn at \(0\)}\label{prop-13-the-dual-of-l-p-spaces-linear-functional-bounded-iff-ctn-at-0}

if \(f:V\rightarrow{\mathbb{K}}\) is a linear functional, TFAE:

\begin{itemize}
\item
  \(f\) bounded
\item
  \(f\) continuous
\item
  \(f\) continuous at \(0 \in V\)
\end{itemize}

\end{proposition}

\begin{proof}

(ii) to (iii): trivial.\\
(i) to (ii): 假设 \(f\) bounded, 那么可以 pick \(C\) s.t. \(\left. |f(v) \middle| \leq C \middle| \middle| v \middle| | \right.\).\\
Pick \(v_{0} \in V,\epsilon > 0\). Set \(\delta := \frac{\epsilon}{C}\). Then

\[\left. | \middle| v - v_{0} \middle| \middle| < \delta\Longrightarrow \middle| f(v) - f(v_{0}) \middle| = \middle| f(v - v_{0}) \middle| \leq C \middle| \middle| v - v_{0} \middle| \middle| < \epsilon \right.\]

从而 ctn.\\
(iii) to (i): \(\exists\delta > 0\) s.t. \(\left. | \middle| v \middle| \middle| \leq \delta\Longrightarrow \middle| f(v) \middle| \leq 1 \right.\).\\
于是 \(\forall v \in V\backslash\left\{ 0 \right\}\), 都有

\[\left. |f(v) \middle| = \frac{|\ f(v \cdot \frac{\delta}{\left. | \middle| v \middle| | \right.})|}{\frac{\delta}{\left. | \middle| v \middle| | \right.}} \leq \frac{\delta}{\left. | \middle| v \middle| | \right.} \right.\]

taking \(C = \frac{1}{\delta}\), 得到 boundedness.

\end{proof}

\begin{qlremark}{Remark}

这个 proposition 看起来很神奇, 把一个整体性质和局部性质等价了, 但是我们知道 linear map 就是局部决定整体的, by its def.\\
recall in 395: 实际上这个性质应该对所有的 linear map 都成立, 不只是 linear functionals.\\
通常我们认为 linear map 总是 ctn 的, 但是其实它 ctn iff bounded, unbounded 的时候就不 ctn.\\
以及: \textbf{linear map between finite dim spaces 总是 bounded 的, 从而总是 ctn 的}. 不过这里我们要讨论的就是 infinite dim spaces. 比如 \(L^{p}\).

\end{qlremark}

% qlnotes: kind=definition; id=def-13-the-dual-of-l-p-spaces-dual-space; concepts=dual-space; aliases=dual space
\begin{definition}{dual space}\label{def-13-the-dual-of-l-p-spaces-dual-space}

If \(V\) is a NVS, 我们定义它的 \textbf{dual space} as:

\[V^{\ast} := \left\{ {\text{bounded linear functionals} f:V\rightarrow{\mathbb{K}}} \right\}\]

\end{definition}

% qlnotes: kind=definition; id=def-13-the-dual-of-l-p-spaces-norm-of-dual-space-dual-norm; concepts=norm-of-dual-space-dual-norm; aliases=norm of dual space: 即 dual norm
\begin{definition}{norm of dual space: 即 \textbf{dual norm}}\label{def-13-the-dual-of-l-p-spaces-norm-of-dual-space-dual-norm}

Given \(f \in V^{\ast}\), set

\[\left. | \middle| f \middle| \middle| {}_{\ast}: = \sup\limits_{v \in V\backslash{\{ 0\}}}\frac{|f(v)|}{\left. | \middle| v \middle| | \right.} = \sup\limits_{||v|| = 1} \middle| f(v)| \right.\]

where \(\parallel v \parallel\) 表示的是 \(V\) 上使用的 norm. 这个 norm 被称为 dual norm.

\end{definition}

这个形式是我们在各种地方见过非常多次的{ }\textbf{} operator norm, 只不过这里, 指定一个 NVS, 对于其 dual space 上的 linear functional, 它是固定的, 不需要指定 \(v\) 和 \(f(v)\)使用哪个 norm, 因为 \(f(v)\) 就是标量, 而 \(v\) from 原 NVS, 已经指定好 norm.

\begin{qlremark}{Remark}

从定义中我们可知, 对于任意的 \(v \in V\), \(f \in V^{\ast}\), 都有:

\[\left. |f(v) \middle| \leq \parallel f\underset{\ast}{\parallel} \middle| v| \right.\]

\end{qlremark}

\subsection{\texorpdfstring{\(V^{\ast}\) being a Banach space}{V\^{}\{\textbackslash ast\} being a Banach space}}\label{vast-being-a-banach-space}

% qlnotes: kind=theorem; id=thm-13-the-dual-of-l-p-spaces-dual-space-is-always-banach; concepts=dual-space-is-always-banach; aliases=dual space is always Banach
\begin{theorem}{dual space is always Banach}\label{thm-13-the-dual-of-l-p-spaces-dual-space-is-always-banach}

对于\textbf{任意的 NVS} \(V\): \(V^{\ast}\) 都是一个 Banach space. (not assuming \(V\) Banach).

\end{theorem}

\begin{proof}

First we can confirm \(V^{\ast}\) is a VS, 因为它由 linear functions of the same size 组成.\\
\textbf{Claim 1: \(V^{\ast}\) 是一个 NVS.}\\
因为任取 \(v \in V,\lambda \in {\mathbb{K}}\) 都有 \(\left. |f(\lambda v) \middle| = \middle| \lambda \middle| \cdot \middle| f(v)| \right.\), 从而

\[\left. f \in V^{\ast},\lambda \in {\mathbb{K}}\Longrightarrow \middle| \middle| \lambda f \middle| \middle| {}_{\ast} = \middle| \lambda \middle| \cdot \middle| \middle| f \middle| |_{\ast} \right.\]

以及

\[\left. f,g \in V^{\ast},v \in V\Longrightarrow \middle| (f + g)(v) \middle| = \middle| f(v) + g(v) \middle| \leq \middle| f(v) \middle| + \middle| g(v)| \right.\]

因而

\[f,g \in V^{\ast}\Longrightarrow \parallel f + g\underset{\ast}{\parallel} \leq \parallel f\underset{\ast}{\parallel} + \parallel g\underset{\ast}{\parallel}\]

下面我们 verify \(V^{\ast}\) Banach.\\
\textbf{Claim 2: 一个 Cauchy seq in \(V^{\ast}\) 一定 pointwise converge to some \(f\).}\\
Pick \((f_{n})_{1}^{\infty}\), 一个 Cauchy seq in \(V^{\ast}\). Let \(\epsilon < 0\), 存在 \(N\) 使得对于任意 \(m,n \geq N\) 都有 \(\parallel f_{n} - f_{m}\underset{\ast}{\parallel} < \epsilon\), 我们简写为:

\[\parallel f_{n} - f_{m}\underset{\ast}{\parallel}\rightarrow 0\]

因而对于任意 \(v \in V\), we have

\[\left. |f_{n}(v) - f_{m}(v) \middle| \leq \parallel f_{n} - f_{m}\underset{\ast}{\parallel} \middle| \middle| v \middle| \middle| \rightarrow 0 \right.\]

并且我们知道 \textbf{\(\mathbb{K}\) 是 complete 的}, 因而 \(f_{n}(v)\) converges in \(\mathbb{K}\) to some element, declared to be \(f(v)\).\\
即 \textbf{\(f_{n}\rightarrow f\) pointwisely}:

\[\lim\limits_{n\rightarrow\infty}f_{n}(v) = f(v)\]

(这是自然的, 因为如果 linear function \(f - g\) 的 operator norm 是 \(0\), 那么说明它们毫无差别, 否则一定有某个地方 \(f,g\) 的 image 不一样, 使得这个 norm 不是 \(0\).)\\
\textbf{Claim 3: \(f\) 是 linear 的, 并且 bounded (从而 ctn), 即 \(f \in V^{\ast}\).}\\
linearity: 由于每个 \(f_{n}\) 都是 linear 的,

\[f_{n}(x + \alpha y) = f_{n}(x) + \alpha f_{n}(y)\]

因而

\[f(x + \alpha y) = \lim\limits_{n\rightarrow\infty}f_{n}(x + \alpha y) = \lim\limits_{n\rightarrow\infty}(f_{n}(x) + \alpha f_{n}(y)) = \lim\limits_{n\rightarrow\infty}f_{n}(x) + \alpha\lim\limits_{n\rightarrow\infty}f_{n}(y) = f(x) + \alpha f(y)\]

因此 \(f\) 是线性的.\\
\textbf{(Note: 这里证明了 linear map 的 pointwise 极限一定也是 linear map.)}\\
Boundedness: Note a standard fact from metric spaces: \textbf{every Cauchy sequence is bounded.}\\
因而 \(f_{n}\) 是一个 bounded seq, 即存在 \(M > 0\) such that \(\parallel f_{n} \parallel \leq M\) for all \(n\). Then

\[\left. |f(x) \middle| = \middle| \lim\limits_{n\rightarrow\infty}f_{n}(x) \middle| \leq \lim\limits_{n\rightarrow\infty} \middle| f_{n}(x) \middle| \leq \lim\limits_{n\rightarrow\infty} \parallel f_{n}\underset{\ast}{\parallel} \parallel x \parallel \leq M\, \parallel x \parallel \right.\]

Hence \(f\) is bounded (continuous), and \(\parallel f\underset{\ast}{\parallel} \leq M\). \textbf{Claim 4: \(\parallel f_{n} - f\underset{\ast}{\parallel}\rightarrow 0\), proving \(V^{\ast}\) 是 Banach 的.} WTS:

\[\left. \parallel f_{n} - f \parallel = \sup\limits_{\parallel x \parallel = 1} \middle| (f_{n} - f)(x) \middle| \rightarrow 0 \right.\]

//TO BE DONE.

\end{proof}

Actually 这个 Theorem 有更 general 的形式:

% qlnotes: kind=theorem; id=thm-13-the-dual-of-l-p-spaces-theorem-002; concepts=theorem-002
\begin{theorem}\label{thm-13-the-dual-of-l-p-spaces-theorem-002}

对于任意 nvm \(V\) 和 Banach \(W\), \(\mathcal{L}(V,W)\) 一定是 Banach 的.

\end{theorem}

Proof 见 Folland 5.4.

\subsection{\texorpdfstring{\((L^{p})^{\ast} = L^{q}\), \(\frac{1}{p} + \frac{1}{q} = 1\)}{(L\^{}\{p\})\^{}\{\textbackslash ast\} = L\^{}\{q\}, \textbackslash frac\{1\}\{p\} + \textbackslash frac\{1\}\{q\} = 1}}\label{lpast-lq-frac1p-frac1q-1}

% qlnotes: kind=theorem; id=thm-13-the-dual-of-l-p-spaces-conjugate-exponent-p-q-l-p-l-q-dual-space; concepts=conjugate-exponent-p-q-l-p-l-q-dual-space; aliases=对于互为 conjugate exponent 的 p,q, L^p 是 L^q 的 dual space
\begin{theorem}{对于互为 conjugate exponent 的 \(p,q\), \(L^{p}\) 是 \(L^{q}\) 的 dual
space}\label{thm-13-the-dual-of-l-p-spaces-conjugate-exponent-p-q-l-p-l-q-dual-space}

For \(1 < p,q < \infty\) with \(\frac{1}{p} + \frac{1}{q} = 1\), we have:

\[(L^{p})^{\ast} = L^{q}\]

In particular the Hilbert space:

\[(L^{2})^{\ast} = L^{2}\]

\end{theorem}

\begin{proof}

Define map

\[\begin{matrix}
L^{q} & {\rightarrow(L^{p})^{\ast}} \\
g & {\mapsto\varnothing_{g}}
\end{matrix}\]

where

\[\varnothing_{g}(f) := \int fg,\quad f \in L^{p}\]

It is well-defined by Hölder:

\[f \in L^{p},g \in L^{q}\Longrightarrow fg \in L^{1}\]

and

\[\left. | \middle| fg \middle| \middle| {}_{1} = \int \middle| fg \middle| \leq \middle| \middle| f \middle| \middle| {}_{p} \middle| \middle| g \middle| |_{q} \right.\]

Easy:

\[\varnothing_{g}(f_{1} + f_{2}) = \varnothing_{g}(f_{1}) + \varnothing_{g}(f_{2})\]

Also

\[\left. |\varnothing_{g}(f) \middle| = \middle| \int fg\  \middle| \leq \int \middle| fg \middle| \leq \middle| \middle| f \middle| \middle| {}_{p} \cdot \middle| \middle| g \middle| |_{q} \right.\]

Thus

\[\varnothing_{g} \in (L^{p})^{\ast}\ \]

\end{proof}

\section{\texorpdfstring{the dual of \(L^{p}\)-II {[}Fol 6.2{]}}{the dual of L\^{}\{p\}-II {[}Fol 6.2{]}}}

\section{\texorpdfstring{the dual of \(L^{p}\)-III {[}Fol 6.2, finished{]}}{the dual of L\^{}\{p\}-III {[}Fol 6.2, finished{]}}}

\printbibliography
\end{document}
